\documentclass[11pt]{article}
\usepackage[a4paper,margin=1in]{geometry}
\usepackage[T1]{fontenc}
\usepackage{lmodern,microtype}
\usepackage{amsmath,amssymb,amsthm,mathtools}
\usepackage{booktabs,array,enumitem,xcolor}
\usepackage[numbers,sort&compress]{natbib}
\usepackage[colorlinks=true,linkcolor=blue!55!black,citecolor=blue!55!black,urlcolor=blue!55!black]{hyperref}
\linespread{1.06}
\newtheorem{theorem}{Theorem}[section]
\newtheorem{proposition}[theorem]{Proposition}
\newtheorem{lemma}[theorem]{Lemma}
\theoremstyle{remark}\newtheorem{remark}[theorem]{Remark}
\newcommand{\C}{\mathbb C}
\newcommand{\card}{\operatorname{card}}

\title{A Rank-Growing Conjugate Resonance-Cloud Atlas\\
\large Shell-Complete Divisors Across Sixteen Noise Scales}
\author{Liang Wang}
\date{July 2026}

\begin{document}
\maketitle

\begin{abstract}
The quartet analysis RH-212--RH-221 ends with an exact obstruction: a fixed
quartic, even if perfectly controlled, cannot provide a growing locally finite
spectral divisor. We therefore construct the first rank-growing cloud on the
same folded Gaussian operators.

At sixteen frozen noise scales
\[
 0.04,\;0.032,\ldots,\;0.00125,
\]
the fine dimension satisfies $N_\sigma\sigma\simeq5.12$ and the second channel
is its Haar compression. The Perron and negative parity resonances are
removed, the remaining spectrum is scaled by the inherited Hardy radius
$0.85$, and a deterministic Arnoldi window is partitioned into real
singletons and nonreal conjugate pairs. A cloud is admitted only as a union
of complete radial shells. The predeclared target ranks are
$k_\ell=4+2\ell$.

All 32 endpoint clouds are conjugate closed. Their actual ranks grow
strictly from four to $34$--$35$ on both channels. Shell completion overshoots
a target by at most one root; at most one nonreal root cut by the inner
Arnoldi-window boundary is discarded. The minimum resolved gap after a
selected cloud is $7.3991\times10^{-5}$ and the recorded conjugacy error is
zero at binary64 precision.

The construction is an operator-derived finite atlas, not an all-level
spectral theorem. In particular the linear rank schedule is a frozen stress
test rather than a canonical asymptotic law. No determinant limit, zeta
identification, or Gate-A closure is asserted.
\end{abstract}

\section{From one quartet to a growing cloud}

RH-219 proves that bounded degree remains bounded under locally uniform
limits, while powers of one quartic grow only multiplicity on fixed support
\citep{WangRH219}. RH-221 consequently designates a rank-growing physical
divisor as the next Gate-A object \citep{WangRH221}. The necessary first task
is deliberately finite: construct larger clouds without breaking the
conjugation symmetry of a real operator.

The point is subtler than taking the first $k$ roots by modulus. A numerical
window can end between the members of a conjugate pair, and a fixed
cardinality prefix can cut a pair after an odd number of real roots has
entered. Either event produces a polynomial with nonreal coefficients and
therefore cannot represent a real spectral factor.

\section{Folded Gaussian endpoints}

Let $f_u(x)=1-ux^2$ at the first band-merging parameter
$u=1.543689012692\ldots$. On positive midpoint nodes
$x_j=(j+1/2)/N$, the folded noisy transition matrix has row weights
\begin{equation}\label{eq:weights}
 W_{ij}=
 \exp\!\left[-\frac{(x_j-f_u(x_i))^2}{2\sigma^2}\right]
 +\exp\!\left[-\frac{(-x_j-f_u(x_i))^2}{2\sigma^2}\right],
 \qquad
 M_{\sigma,N}(i,j)=\frac{W_{ij}}{\sum_m W_{im}}.
\end{equation}
Only entries within eight Gaussian standard deviations are retained, after
which each row is renormalized. The default cutoff is below binary64
roundoff at the level of omitted Gaussian mass.

For every frozen $\sigma$, set
\begin{equation}
 N_\sigma=
 \max\!\left(32,\,
 2\,\operatorname{round}\frac{5.12}{2\sigma}\right).
\end{equation}
The fine endpoint is $M_{\sigma,N_\sigma}$. If
$E:\C^{N_\sigma/2}\to\C^{N_\sigma}$ is the pair-average Haar isometry,
\begin{equation}
 E_{2j,j}=E_{2j+1,j}=2^{-1/2},
\end{equation}
the right endpoint is $E^*M_{\sigma,N_\sigma}E$. Both matrices are real.

The historical Hardy normalization uses $r_H=0.85$. Resolve a leading
eigenvalue window, remove the root nearest $1$ and the most negative resolved
real root, and divide every remaining value by $r_H$. These are the bulk
resonances used below. Scaling does not affect conjugation or radial order.

\section{Conjugate shells}

\begin{lemma}[Real spectral parity]\label{lem:conj}
The eigenvalue multiset of a real finite matrix is invariant under complex
conjugation, including algebraic multiplicity.
\end{lemma}
\begin{proof}
Its characteristic polynomial has real coefficients. Hence
$p(\overline z)=\overline{p(z)}$, and root multiplicities are preserved.
\end{proof}

A \emph{shell} is either a resolved real root $\{\lambda\}$ or a nonreal pair
$\{\lambda,\overline\lambda\}$. Shells are ordered by decreasing outer
modulus. Given target $k$, include whole shells until their cumulative
cardinality first reaches $k$.

\begin{proposition}[Shell-complete prefix]\label{prop:prefix}
Every selected cloud is conjugate closed and is the zero multiset of a monic
real polynomial. Its rank is either $k$ or $k+1$.
\end{proposition}
\begin{proof}
Each shell is conjugate closed. A union of shells is therefore conjugate
closed, so its monic root polynomial has real coefficients by
Lemma~\ref{lem:conj}. Every shell has size one or two. Immediately before
the last shell, cumulative size is below $k$; adding at most two roots gives
rank at most $k+1$.
\end{proof}

The proposition does not require simple eigenvalues. Repeated real roots can
be represented by singleton copies, and a repeated nonreal root carries the
same number of conjugate copies.

\subsection{Candidate-window boundary}

Arnoldi is asked for $k+16$ eigenvalues before Perron/parity removal. Its
returned inner boundary may itself cut a nonreal pair. Such an unmatched
boundary root is discarded; no conjugate value is synthesized. Selection
continues only if the remaining complete shells still reach the target.
The archive records the discarded count. It is at most one in every endpoint.

This rule separates two claims:
\begin{enumerate}[leftmargin=2em]
\item selected clouds are exactly conjugate closed within numerical tolerance;
\item the finite Arnoldi window is not claimed to resolve the full infinite
spectrum.
\end{enumerate}

\section{Frozen rank ladder}

The sixteen targets are
\begin{equation}\label{eq:ranks}
 k_\ell=4+2\ell,\qquad 0\le\ell\le15.
\end{equation}
This schedule was fixed before inspecting the resulting lower shells. It
forces a genuine degree increase while keeping a sixteen-root candidate
margin. It is not derived from a Weyl law and is not promoted to a canonical
rank/noise relation.

Representative endpoint data are:
\begin{center}
\begin{tabular}{rrrrrr}
\toprule
$\sigma$ & target & left rank & right rank &
left inner modulus & right inner modulus\\
\midrule
$0.04000$ & 4  & 4  & 4  & $0.35946$ & $0.35347$\\
$0.02500$ & 8  & 9  & 9  & $0.19658$ & $0.19027$\\
$0.01000$ & 16 & 16 & 17 & $0.13873$ & $0.12899$\\
$0.00500$ & 22 & 23 & 23 & $0.13870$ & $0.13553$\\
$0.00250$ & 28 & 29 & 29 & $0.07384$ & $0.08025$\\
$0.00125$ & 34 & 35 & 34 & $0.08177$ & $0.10716$\\
\bottomrule
\end{tabular}
\end{center}

Actual ranks are strictly increasing in scale index on each channel. The
frequent odd ranks are not defects: they indicate that one or more real bulk
resonances lie before the closing conjugate shell.

\section{Radial isolation and global gauge}

If $\Lambda$ is the selected cloud and $\Lambda^+$ is the first omitted
complete shell, define
\begin{equation}
 g(\Lambda)=\min_{\lambda\in\Lambda}|\lambda|
 -\max_{\mu\in\Lambda^+}|\mu|.
\end{equation}
All 32 resolved gaps are positive; the minimum is
$7.3991\times10^{-5}$. This is a finite separation certificate inside the
candidate window, not a lower bound uniform in $\sigma$.

For later work the archive also records one global affine gauge:
\begin{equation}\label{eq:gauge}
 m_\Lambda=\frac1{\card\Lambda}\sum_{\lambda\in\Lambda}\lambda,\qquad
 s_\Lambda^2=\frac1{\card\Lambda}
 \sum_{\lambda\in\Lambda}|\lambda-m_\Lambda|^2,
 \qquad q_\lambda=\frac{\lambda-m_\Lambda}{s_\Lambda}.
\end{equation}
Conjugation makes $m_\Lambda$ real up to roundoff. The left/right center
difference is at most $0.01610$ and the RMS-radius difference at most
$0.01490$ in the frozen atlas. These are diagnostics, not convergence rates.

Reciprocal points $1/\lambda$ are stored at the same time. They will matter
because Fredholm determinants have reciprocal resonance zeros; the present
paper does not yet use that fact.

\section{Reproducibility controls}

The eigensolver receives a deterministic trigonometric start vector. Every
endpoint stores the full resolved candidate list, complete-shell count,
selected roots, normalized roots, reciprocal roots, Perron/parity values,
radial gap, global gauge, and a whole-matrix Frobenius ledger. Unit tests
check:
\begin{enumerate}[leftmargin=2em]
\item an artificial split pair is repaired by minimal shell completion;
\item an unmatched inner-boundary root is discarded;
\item global centering and RMS normalization satisfy their identities;
\item reciprocal multiplication returns one.
\end{enumerate}

The archived result contains 32 endpoint records and sixteen channel
comparisons. It can therefore be reused without rerunning Arnoldi or changing
the root-selection rule.

\section{Claim boundary and next question}

A rank-growing cloud now exists at every frozen endpoint, so the fixed-degree
wall has been passed at the finite numerical level. Three stronger statements
remain open:
\begin{enumerate}[leftmargin=2em]
\item certification of the infinite-spectrum ordering and gaps;
\item a canonical relation between physical scale and selected rank;
\item a locally uniform determinant or locally finite limiting divisor.
\end{enumerate}

The immediate next layer tests whether the shell-complete selection is stable
under candidate-window changes and quantifies how often a naive fixed-rank
prefix violates conjugation. Gates A--E remain open. No self-adjoint
operator, arithmetic trace formula, or Riemann-hypothesis consequence is
claimed.

\bibliographystyle{plainnat}
\bibliography{references}
\end{document}
